\section{Implementation and Safety Engineering}
\label{sec:implementation}

\subsection{Three-tier no\_std policy}

The workspace separates three environments:

\begin{enumerate}
  \item \textbf{Core and VM:} \code{no_std}, no allocation, no unsafe. The
        compiler, lifter, IR, MBA arena, ISA derivation, cipher, payload
        builder, and VM engine use fixed-capacity data structures and only
        \code{core}.
  \item \textbf{Binary formats:} \code{no_std} plus \code{alloc}, no unsafe.
        Parsing and editing ELF/PE images require variable-size buffers in the
        CLI context, but the parsers do not use raw pointers.
  \item \textbf{Protector CLI:} \code{std}, safe Rust. Filesystem access,
        process spawning, temporary directories, and error reporting are
        confined to this boundary. The generated runtime is not linked against
        \code{std}.
\end{enumerate}

The \code{sirus-hardened} crate is the only shipped crate that permits unsafe
code, and it is opt-in. A source gate scans the workspace for unsafe constructs
outside that crate.

\subsection{Fixed capacities}

All compiler and runtime collections are fixed-capacity. The main bounds are
shown in Table~\ref{tab:limits}. A capacity overflow is a typed build error, not
a panic or an allocation.

\begin{table}[htbp]
\centering
\small
\caption{Fixed capacities in the compiler and runtime.}
\label{tab:limits}
\begin{tabular}{@{}ll@{}}
\toprule
Limit & Value \\
\midrule
Functions per protection run & 8 \\
Blocks per function & 256 \\
IR operations per function & 8192 \\
Function size & 16 KiB \\
Payload size & 128 KiB \\
Private read-only data & 2 KiB \\
Private stack & 4 KiB \\
Virtual registers & 96 (GPRs, flags, constants, scratch, XMM halves) \\
Manifest entries & 256 \\
Constant-time passes & 1--4096 (default 256) \\
Default fuel & 2,000,000 steps \\
\bottomrule
\end{tabular}
\end{table}

The capacities are chosen so that a protected function and its VM fit in a
single RX segment and in the test-thread stack used by the differential harness.
They are not arbitrary: the 128 KiB payload ceiling was raised after a
six-function max-profile differential build exceeded the previous 64 KiB limit,
and it remains small enough to keep the VM stack frame bounded.

\subsection{Panic elimination}

Normal execution has no panic path. Arithmetic is wrapping, shifts are masked,
division is absent, indexing is checked through slice accessors, and every
malformed state returns a status. The generated panic handler exists only as a
final \code{no_std} link guard; it is not reachable from a supported semantic
operation.

The test-only lint exemptions allow assertion panics in integration tests. The
shipped code keeps the ban on \code{unwrap}, \code{expect}, \code{panic},
\code{todo}, \code{unimplemented}, and \code{unreachable}. Clippy runs with
\code{-D warnings} across all targets; the release gate includes rustdoc with
denied warnings and a dependency policy check.

\subsection{Determinism and reproducibility}

A protected image is a deterministic function of the input bytes, the seed, the
profile, and the pinned toolchain. The seed drives every random choice through a
single build RNG; no OS entropy is mixed after the seed is fixed. The
reproducible-output check protects normal, hardened, and constant-time builds and
compares their hashes.

Omit \verb|--seed| to derive a seed from OS entropy. A fixed seed is useful for
reproducible builds, regression tests, and differential comparison. The
protector does not embed the seed in the output.

The repository pins stable Rust 1.99.0 and the Windows GNU target in
\path{rust-toolchain.toml}. The generated VM is compiled by \code{rustc} at
protection time with the flags in Section~\ref{sec:runtime}. The protector uses
\code{cc}, LLD, and \code{objcopy} on Linux, and
\code{x86_64-w64-mingw32-gcc} plus \code{lld-link} on Windows.

\subsection{Compiler/runtime source sharing}

The protector embeds the VM engine, the shared ISA constants, the cipher, and
the derivation module with \code{include_str!}. A build script tracks every
embedded file and triggers a rebuild when one changes. This addresses a subtle
class of release bugs: a stale embedded engine can disagree with the compiler's
wire format even though the source tree looks consistent. The release history
contains one such bug (a missing semantic import that became a catch-all match
arm in release builds), which is now covered by a release-mode test run and the
build-script dependency tracking.

\subsection{Threading, reentrancy, and signals}

The VM uses only per-invocation stack state. There is no mutable global engine
state and no shared mutable cache. Separate host threads can call protected
functions concurrently without sharing virtual registers or private memory. An
eight-thread fixture calls one protected function 50,000 times per thread and
compares the aggregate result with native.

A signal handler can re-enter a protected function while the main thread is
inside it. The generated entry uses a fresh stack frame and the Linux shim saves
and restores the host context, so the interrupted and reentrant calls are
independent. A SIGALRM fixture runs 200 protected calls while a 50 ms interval
timer re-enters the same function from the handler; both the interrupted result
and the handler result match native, and the handler must fire at least once.
This test exposed and now covers a real stack-alignment bug in an earlier shim.

\subsection{Logging and diagnostics}

The generated runtime has no logging, formatting, or string constants for
diagnostics. The optional \code{trace} feature of the \code{sirus-vm} reference
crate prints constant-time events for debugging; it is not enabled in generated
artifacts and its default sink is a no-op. This keeps the injected segment free
of diagnostic strings and avoids a hidden stderr dependency.

\subsection{Code and artifact size}

The default benchmark build adds about 30 KiB of VM and payload to a 4.5 MiB
binary (0.7\%). The payload for a single benchmark function is about 13.4 KiB;
the VM segment including the engine and tables is larger. The exact size depends
on the function, profile, and seed. The injected segment is read-execute and
non-writable; the host image is not compressed or packed.
